\section{Introduction}

Tool-using agents already behave like institutional actors: they receive delegated authority, buy services, call tools hosted by other organizations, emit records that become accounting evidence, and trigger irreversible actions. Deployments can record who configured a prompt or approved an integration, but cannot prove that a later cross-vendor action satisfied a declared constitution, treaty, and capability boundary before crossing the trust boundary. We close that gap by reducing institutional authority to four decidable operations on a polity triple $(T, C, K)$: admission against the receipt scope $T$, attenuation within the citizenship roster $C$, bilateral treaty intersection between two constitutions $K_1$ and $K_2$, and amendment refinement within $K$. The implementation checks these operations at the Rust kernel boundary and signs admitted decisions into canonical receipts; the formal layer is narrower, with two load-bearing Lean theorems for treaty admission and ladder-floor stability, two definitional bridges for amendment refinement, and explicit runtime assumptions left outside the mechanized proof.

The construction rests on three substrate conditions. Authority is carried by attenuating capability tokens, not ambient credentials. Every kernel decision is recorded as a signed receipt over canonical JSON bytes. Membership, treaty, and amendment rules are predicates whose safety claims map to machine-checked proof artifacts. Chio implements that stack: an untrusted agent presents a capability, a trusted runtime kernel evaluates scope and guard predicates, tool servers remain isolated, a capability authority issues and revokes bounded grants, and the receipt log records the decision.

The polity claim is declarative. A Chio polity has final authority over its receipt-admission boundary; it has none over physical persons, land, or external courts. Its territory is a receipt namespace, its citizens are Decentralized Identifiers (DIDs) per W3C DID Core~\cite{w3cDidCore} or kernels admitted into that namespace, its border is the runtime admission hook, and its courts are verdict artifacts over signed receipts. The substrate gives a constructive test for any further claim: it lives or fails as an admission predicate, treaty artifact, proof term, or replayable script. Section~\ref{sec:discussion} positions this against Montevideo and network-state literature; Section~\ref{sec:limits} states the legal scope.

The contributions are five falsifiable artifacts.

\begin{itemize}[noitemsep,topsep=2pt,leftmargin=1.2em]
  \item A formal model relating polity scope, constitutional predicates, receipt admission, treaty intersection, and amendment refinement (Section~\ref{sec:model}).
  \item Two load-bearing Lean theorems and two definitional refinement bridges in \codepath{formal/lean4/Chio/Chio/Treaty/Intersection.lean}: \thm{treaty_admission_iff_predicate_intersection} and \thm{treaty_admission_stable_under_ladder_floor} (load-bearing), with \thm{amendment_admissible_iff_backward_refinement} and \thm{amendment_without_refinement_rejected} (definitional, anchoring the type-level enactment invariant) (Section~\ref{sec:model}).
  \item An implementation of the model in the Chio runtime, covering treaty scope and ladder intersection, pre-dispatch admission, strict bilateral DSSE verification, and a multi-lane anchor whose maturity varies by lane (Section~\ref{sec:implementation}).
  \item A three-vendor buyer-closure demonstration built on a fixture-backed generator at \codepath{examples/chiodos-3vendor/}, exercising every load-bearing predicate in the substrate (Sections~\ref{sec:implementation} and \ref{sec:evaluation}).
  \item An evaluation that measures dispatch admission, treaty intersection, selective disclosure, and replay-corpus stability, and that names which lanes and primitives have no current measurement (Section~\ref{sec:evaluation}).
\end{itemize}
